\documentclass[10pt,aspectratio=169]{beamer}

% ---------- Theme ----------
\usetheme{Madrid}
\usecolortheme{default}
\setbeamertemplate{navigation symbols}{}

% ---------- Packages ----------
\usepackage{booktabs}
\usepackage{amsmath, amsfonts, amssymb, bm}
\usepackage{graphicx}
\usepackage{array}
\usepackage{multirow}
\usepackage{siunitx}

% ---------- Notation macros (edit these to match the paper) ----------
\newcommand{\drug}{v}          % drug consumption
\newcommand{\nondrug}{c}       % non-addictive consumption
\newcommand{\house}{h}         % housing services
\newcommand{\leis}{\ell}       % leisure
\newcommand{\uutil}{\nu}       % per-period utility function symbol
\newcommand{\disc}{\beta}      % discount factor

% ---------- Title info ----------
\title[Housing First vs Treatment First]{Ending Homelessness: Housing First or Treatment First?}
\author[Ay\c{s}e \.{I}mrohoroglu, Shahar Rotberg, Kai Zhao]{Ay\c{s}e \.{I}mrohoroglu \and Shahar Rotberg \and Kai Zhao}
\institute[USC/CMHC/UConn]{USC \quad\quad CMHC \quad\quad UConn}
\date{}

\begin{document}

% ---------------- Title ----------------
\begin{frame}
  \titlepage
\end{frame}

% ---------------- Intro ----------------
\begin{frame}{Homelessness and addiction in United States}
\begin{itemize}
\item Homelessness is a pressing issue in the U.S.
\item 568{,}000 people were homeless on a given night in 2019 in U.S.
\item 1{,}416{,}908 experienced sheltered homelessness throughout 2017
\item 4\% in general U.S. population use drugs; 2\% serious, 2\% light (SAMHSA 2023)
\item 25--50\% of homeless have a substance use disorder (SUD) (Burt et al., 1999)
\item SUDs contribute substantially to morbidity and mortality (Coombs et al., 2024)
\end{itemize}
\end{frame}

\begin{frame}{Fighting homelessness}
\begin{itemize}
\item Addiction and homelessness are tightly connected; addressing one affects the other.
\item Main tools to fight homelessness:
\begin{itemize}
\item Currently: \emph{Housing First} (HF)
\item Alternative: \emph{Treatment First} (TF)
\end{itemize}
\end{itemize}
\end{frame}

\begin{frame}{Decades-long policy debate}
\begin{itemize}
\item TF advocates argue that HF creates substantial moral hazard.
\item Thus, HF leads to more addiction and thus more homelessness.
\item HF camp says moral hazard is low and homelessness hinders sobriety and employment.
\item TF would lead to more addiction and more homelessness.
\end{itemize}
\end{frame}

\begin{frame}{This paper}
\begin{itemize}
\item Develop a quantitative economic model with endogenous homelessness and addiction.
\item Compare HF with TF policies:
\begin{itemize}
\item HF: No conditions on drug use; receive housing vouchers with a probability.
\item TF: Must be drug-free to have a chance to receive housing vouchers.
\end{itemize}
\end{itemize}
\end{frame}

\begin{frame}{Preview of results}
\begin{itemize}
\item TF increases homelessness immediately, but reduces addiction in the long run.
\item But HF generates higher welfare both in short run and in long run.
\item Direction holds for low moral hazard (benchmark) and high moral hazard.
\end{itemize}

\vspace{0.5em}
\begin{table}
\centering
\small
\begin{tabular}{lccc}
\toprule
Statistic & Short run, PE & Long run, PE & Long run, GE\\
\midrule
Homelessness & $\Uparrow$ & $\Uparrow$ & $\Uparrow$\\
Chronic homelessness & $\Uparrow$ & $\Uparrow$ & $\Uparrow$\\
Addiction & $\Leftrightarrow$ & $\Downarrow$ & $\Downarrow$\\
Newborns' welfare & $\Downarrow$ & $\Downarrow$ & $\Downarrow$\\
\bottomrule
\end{tabular}
\end{table}
\end{frame}

% ---------------- Model ----------------
\section{Model}

\begin{frame}{High-level model overview}
\centering
\includegraphics[width=0.92\textwidth]{slide_imgs/page_08.png}
\end{frame}

\begin{frame}{Model --- General features}
\begin{itemize}
\item Overlapping generations with heterogeneous individuals and incomplete markets.
\item Homelessness as in \.{I}mrohoroglu and Zhao (2024).
\item Rational drug use inspired by Becker and Murphy (1988).
\item Probabilistic addiction as in Greenwood et al. (2024/2025).
\item Individuals value $\nondrug$, $\drug$, $\house$, and leisure $\leis$ and face risks: death, labor income, addiction.
\end{itemize}
\end{frame}

\begin{frame}{Model --- Demographics}
\begin{itemize}
\item Individuals live from age $1$ to age $J$, retire at age $J_R$.
\item Survive from $j$ to $j+1$ with probability $\phi(j)$; die with certainty after age $J$.
\item Deceased net assets are equally distributed to everyone (bequests), denoted $\pi_a$.
\end{itemize}
\end{frame}

\begin{frame}{Model --- Income}
\begin{itemize}
\item Labor income depends on:
\begin{itemize}
\item Fixed component $x$ (lifetime), drawn from $\Omega_x(x)$.
\item Idiosyncratic component $z$, Markov chain $\Omega_z(z,z')$.
\item Age productivity profile $\{\zeta_j\}_{j=1}^{J_R-1}$.
\end{itemize}
\item Employment status $e\in\{0,1\}$ follows Markov chain $\Omega_e(\cdot,\cdot)$.
\item Unemployed may receive UI with probability $p^{ui}$:
\[
y^{ui}_j =
\begin{cases}
\min\{ v^e(1-\leis) x z \zeta_j,\ \bar y^{\leis}\} & \text{if UI}\\
0 & \text{otherwise.}
\end{cases}
\]
\end{itemize}
\end{frame}

\begin{frame}{Model --- Income continued}
\begin{itemize}
\item Retirees receive Social Security $y^R(x,z)=\kappa(x,z)\bar y^{\leis}$, based on $x,z$ at retirement.
\item Define income $\hat y^{\leis}_j$:
\[
\hat y^{\leis}_j=
\begin{cases}
y^{\leis}_j & e=0\\
y^{ui}_j & e=1\\
y^R & j\ge J_R.
\end{cases}
\]
\item Capital income: $\hat y^{k}_j=r a$ (if $a<0$, this is mortgage interest payments).
\end{itemize}
\end{frame}

\begin{frame}{Model --- Housing and mortgage markets}
\begin{itemize}
\item Individuals can rent or own ($o\in\{0,1\}$) or be homeless ($\house=\underline{h}$).
\item Renters choose $h'\in\{h_i\}_{i=1}^n$ and pay per-unit rent $p^r$.
\item Homeowners choose $h'\in\{h_i\}_{i=n-1}^n$ and pay per-unit price $p^h$.
\item One-period mortgages at interest rate $r$, with LTV constraint $\lambda p^h h \ge -a$ (benchmark $\lambda=0.8$).
\item Selling cost $\tau_s=0.06$.
\end{itemize}
\end{frame}

\begin{frame}{Model --- Homelessness}
\begin{itemize}
\item Homeless receive reduced housing services $\underline{h}<h_1$ at no cost.
\item Higher mortality: multiplier $\kappa_{h,\phi}=1.6$ on baseline death risk.
\item Worse labor-market transitions:
\begin{itemize}
\item Re-employment multiplier $\kappa_{h,ue}=0.8$.
\item Unemployment-risk multiplier $\kappa_{h,eu}=1.25$.
\end{itemize}
\item Lower probability of addiction recovery when homeless (multiplier $\kappa_{H,d}=\tfrac{1}{6}$).
\end{itemize}
\end{frame}

\begin{frame}{Model --- Housing vouchers}
\begin{itemize}
\item Low-income renters eligible for a voucher lottery if:
\begin{itemize}
\item Gross income below 50\% of median ($\kappa_V=0.5$), and
\item Required rent contribution rate is 30\% of income ($\kappa_r=0.3$).
\end{itemize}
\item Voucher covers the gap between rent and $\kappa_r$ of income (subject to program rules).
\item Benchmark lottery intensity: per-period receipt probability among eligible individuals $\pi^{V}=0.06$,
which implies about 26.6\% of eligible receive vouchers over a year (calibration target).
\end{itemize}
\end{frame}

\begin{frame}{Model --- Food stamps}
\begin{itemize}
\item Low-income individuals receive food stamps with certainty.
\item Food stamps must be spent on non-addictive consumption: $\nondrug \ge \omega^{food}$.
\end{itemize}
\end{frame}

\begin{frame}{Model --- Drug experimentation and addiction}
\begin{itemize}
\item Individuals may experiment ($q\in\{0,1\}$) with $\drug$ until age $J_A$.
\item Experimenters become addicted ($d\in\{0,1\}$) with probability $\Omega_d(0,1\mid q=1)$.
\item Heavy experimenters face higher addiction risk: $\pi_{d,high}=6.48\times\pi_{d,low}$ (NSDUH 2024).
\item Addicts have:
\begin{itemize}
\item Higher mortality: multiplier $\kappa_{d,\phi}=6.9$ while addicted.
\item Worse labor-market transitions: re-employment multiplier $\kappa_{d,ue}=0.37$ and unemployment-risk multiplier $\kappa_{d,eu}=2.7$.
\item Recovery probability each period if housed: $\Omega_d(1,0\mid \house>\underline{h})=0.025$.
\item Recovery probability if homeless: $\Omega_d(1,0\mid \house=\underline{h})=\tfrac{1}{6}\times 0.025$.
\end{itemize}
\end{itemize}
\end{frame}

\begin{frame}{Model --- Income taxes}
\begin{itemize}
\item Progressive labor income taxes (Heathcote et al., 2014):
\[
\rho^{\leis}_j = \hat y^{\leis}_j - \tau_{\leis}\,(\hat y^{\leis}_j)^{\xi}.
\]
\item $\tau_{\leis}$ pins down average labor income taxes; $\xi$ controls progressivity.
\item Positive capital income taxed at constant rate $\tau_k$:
\[
\rho^k_j=\tau_k \max\{\hat y^k_j,0\}.
\]
\end{itemize}
\end{frame}

\begin{frame}{Model --- Preferences (non-addicts)}
\small
Non-addicts ($d=0$):
\[
\uutil(\nondrug,\house',\drug,\leis\mid d=0)=
\frac{\Big[(\nondrug^{1-\theta_1}(\house')^{\theta_1})^{1-\theta_{0,2}}(\drug+b)^{\theta_{0,2}}\Big]^{1-\sigma}}{1-\sigma}
+\omega_{\leis}\frac{\leis^{1-\sigma_{\leis}}}{1-\sigma_{\leis}}.
\]
\begin{itemize}
\item $\theta_1(1-\theta_{0,2})$: weight on housing $\house'$.
\item $(1-\theta_1)(1-\theta_{0,2})$: weight on non-drug consumption $\nondrug$.
\item $\theta_{0,2}$: weight on drugs; relative price $p^{v}$.
\item $b$ controls temptation (if $b>0$, not all individuals use drugs).
\end{itemize}
\end{frame}

\begin{frame}{Model --- Preferences (addicts)}
\small
Addicts ($d=1$):
\[
\uutil(\nondrug,\house',\drug,\leis\mid d=1)=
\frac{\Big[(\nondrug^{1-\theta_1}(\house')^{\theta_1})^{1-\theta_{1,2}}(\drug+b)^{\theta_{1,2}}\Big]^{1-\sigma}}{1-\sigma}
+\omega_{\leis}\frac{\leis^{1-\sigma_{\leis}}}{1-\sigma_{\leis}}.
\]
\begin{itemize}
\item $\theta_{1,2}>\theta_{0,2}$: addiction raises the marginal utility/weight of drug consumption (Becker \& Murphy, 1988).
\end{itemize}
\end{frame}

\begin{frame}{Model --- Budget constraints (homeowners)}
\small
\begin{itemize}
\item Choosing to own ($o'=1$) starting as homeowner ($o=1$):
\[
\nondrug + p^v \drug + p^h h(\mathbb{I}_{\{h'\neq \underline{h}\}}\tau_s + \delta_h) + p^h h' + a' \\
= \big(\hat y^{\leis}_j-\rho^{\leis}_j\big) + \big(\hat y^k_j-\rho^k_j\big) + a + p^h h + \omega^{food} + \pi_a.
\]
\item Choosing to own ($o'=1$) starting renter/homeless ($o=0$):
\[
\nondrug + p^v \drug + p^h h' + a'
= \big(\hat y^{\leis}_j-\rho^{\leis}_j\big) + \big(\hat y^k_j-\rho^k_j\big) + a + \omega^{food} + \pi_a.
\]
\end{itemize}
\end{frame}

\begin{frame}{Model --- Budget constraints (renters)}
\small
\begin{itemize}
\item Renting ($o'=0$, $h'>\underline{h}$) starting homeowner ($o=1$):
\[
\nondrug + p^v \drug + p^r h' + p^h h(\delta_h+\tau_s) + a' \\
= \big(\hat y^{\leis}_j-\rho^{\leis}_j\big)+\big(\hat y^k_j-\rho^k_j\big)+a+p^h h+\omega^{food}+\eta\omega^{voucher}+\pi_a.
\]
\item Renting ($o'=0$, $h'>\underline{h}$) starting renter/homeless ($o=0$):
\[
\nondrug + p^v \drug + p^r h' + a' \\
= \big(\hat y^{\leis}_j-\rho^{\leis}_j\big)+\big(\hat y^k_j-\rho^k_j\big)+a+\omega^{food}+\eta\omega^{voucher}+\pi_a.
\]
\end{itemize}
\end{frame}

\begin{frame}{Model --- Budget constraints (homeless)}
\small
\begin{itemize}
\item Choosing homelessness ($o'=0$, $h'=\underline{h}$) starting renter ($o=0$, $h>\underline{h}$):
\[
\nondrug + p^v \drug + a'
= \big(\hat y^{\leis}_j-\rho^{\leis}_j\big)+\big(\hat y^k_j-\rho^k_j\big)+a+\omega^{food}+\pi_a.
\]
\item Choosing homelessness ($o'=0$, $h'=\underline{h}$) starting homeowner ($o=1$):
\[
\nondrug + p^v \drug + p^h h(\delta_h+\tau_s) + a' \\
= \big(\hat y^{\leis}_j-\rho^{\leis}_j\big)+\big(\hat y^k_j-\rho^{\leis}_j\big)+a+p^h h+\omega^{food}+\pi_a.
\]
\end{itemize}
\end{frame}

\begin{frame}{Model --- Individual optimization}
\small
State: $g=(a,x,z,e,\omega_e,h,o,\eta,d)$.
\[
V_j(g)=\max_{\{\leis,a',\nondrug,o',h',\drug,q\}}
\left\{\mathcal{N}+\uutil(\nondrug,h',\drug,\leis\mid d)+
\disc\,\phi(j\mid h',d)\,V_{j+1}(g'\mid z,e,\omega_e,h',q,d)\right\}
\]
subject to budget constraints, LTV, food-stamps, voucher, and housing constraints.

\vspace{0.25em}
Policy functions: $\leis_j(g)$, $a'_j(g)$, $\nondrug_j(g)$, $o'_j(g)$, $h'_j(g)$, $\drug_j(g)$, $q_j(g)$.

\vspace{0.25em}
$\mathcal{N}$ is set so everyone has a positive value of life (average of \$10M).
\end{frame}

\begin{frame}{Housing and rental supply}
\begin{itemize}
\item Competitive construction: $I_h=\epsilon_1 (p^h)^{\epsilon_2}$.
\item Housing law of motion: $H'=(1-\delta_h)H + I_h$.
\item Competitive rental company chooses rental stock $X$ and faces convex management costs:
\[
C_r(X)=\frac{\gamma_1}{\gamma_2}X^{\gamma_2}.
\]
\item Rental firm's dynamic problem:
\[
W(X)=\max_{X'}\left\{p^rX' - C_r(X') - p^h(X'-X) - p^h\delta_h X + \frac{1}{1+r}W(X')\right\}.
\]
\item Optimal rent:
\[
p^r=\gamma_1 (X')^{\gamma_2-1}+\left(\frac{r+\delta_h}{1+r}\right)p^h.
\]
\end{itemize}
\end{frame}

\begin{frame}{Government}
\begin{itemize}
\item Government holds total voucher spending fixed across simulations.
\item Let $\psi_j(g)$ be the distribution of age-$j$ individuals across state space.
\item Budget condition (vouchers):
\[
\sum_{j=1}^{J}\int \eta\,\omega^{voucher}\,\psi_j(g)\,dg = \text{benchmark voucher spending}.
\]
\end{itemize}
\end{frame}

\begin{frame}{Equilibrium definition}
\begin{itemize}
\item Equilibrium: prices, value/policy functions, and distributions such that:
\begin{itemize}
\item Household policy functions solve the DP.
\item Rental firm's policy solves its problem.
\item Government budget clears.
\item Bequests equal total net assets of deceased.
\item Distributions are consistent with optimal choices.
\item Housing and rental markets clear.
\end{itemize}
\end{itemize}
\end{frame}

% ---------------- Calibration ----------------
\section{Calibration}

\begin{frame}{External calibration --- Demographics and preferences}
\small
\begin{table}
\centering
\begin{tabular}{lll}
\toprule
Parameter & Description & Value / Source\\
\midrule
$J$ & Lifespan (model periods) & 60 (ages 26--85)\\
$J_R$ & Retirement age (model periods) & 41 (age 66)\\
$\{\phi^{base}(j)\}_{j=1}^J$ & Baseline survival probs. & Arias (2014)\\
$\sigma$ & Risk aversion & 2\\
$\sigma_{\ell}$ & Leisure curvature & 3 (Chetty 2012)\\
\bottomrule
\end{tabular}
\end{table}
\end{frame}

\begin{frame}{External calibration --- Labor income, unemployment, retirement}
\small
\begin{table}
\centering
\begin{tabular}{lll}
\toprule
Parameter & Description & Value / Source\\
\midrule
$\{\zeta_j\}$ & Life-cycle productivity profile & Guvenen et al. (2023)\\
$\rho_z$ & Persistence of $z$ & 0.9 (Guvenen et al. 2023)\\
$\sigma_z$ & Volatility of $z$ & 0.2 (Guvenen et al. 2023)\\
$\rho_x$ & Intergenerational persistence of $x$ & 0.5 (Guvenen et al. 2023)\\
$\sigma_x$ & Innovation vol. of $x$ & 0.309 (Guvenen et al. 2023)\\
$y^R(x,z)$ & Social Security benefits & Guvenen et al. (2023)\\
$\pi^{\kappa}$ & UI receipt prob. if unemployed & 0.3 (FRED 2014--2024)\\
$\kappa^{ui}$ & UI replacement rate & 0.5\\
$\kappa_{h,ue}$ & Re-employment multiplier if homeless & 0.8 (Desmond \& Gershenson 2016)\\
$\kappa_{h,eu}$ & Unemployment-risk multiplier if homeless & 1.25 (Desmond \& Gershenson 2016)\\
\bottomrule
\end{tabular}
\end{table}
\end{frame}

\begin{frame}{External calibration --- Interest rates and taxes}
\small
\begin{table}
\centering
\begin{tabular}{lll}
\toprule
Parameter & Description & Value / Source\\
\midrule
$r$ & Real interest rate & 3\% (Kaplan et al. 2020)\\
$\tau_k$ & Capital income tax rate & 15\% (Karlman et al. 2021)\\
$\xi$ & Progressivity of labor taxes & 0.85 (Heathcote et al. 2017)\\
\bottomrule
\end{tabular}
\end{table}
\end{frame}

\begin{frame}{External calibration --- Homelessness and addiction}
\small
\begin{table}
\centering
\begin{tabular}{lll}
\toprule
Parameter & Description & Value / Source\\
\midrule
$p^v$ & Price of drugs & 1 (normalization)\\
$\kappa_{h,\phi}$ & Mortality multiplier if homeless & 1.6 (Meyer et al. 2025)\\
$\kappa_{H,d}$ & Recovery prob. if housed & 0.025 (Vaillant 1995)\\
$\pi_{d,high}$ & Addiction risk (heavy) & $6.48\times\pi_{d,low}$ (NSDUH 2024)\\
$\chi_H$ & Gov. cost per homeless & $0.252\,y$ (NAEH 2016)\\
$\chi_D$ & Gov. cost per addict & $0.212\,y$ (CEA 2017)\\
\bottomrule
\end{tabular}
\end{table}
\end{frame}

\begin{frame}{External calibration --- Housing and housing finance}
\small
\begin{table}
\centering
\begin{tabular}{lll}
\toprule
Parameter & Description & Value / Source\\
\midrule
$\delta$ & Housing depreciation rate & 0.6\% (Rosenthal 2014)\\
$\tau_s$ & Selling cost & 6\% (Gruber \& Martin 2003)\\
$\varsigma_2$ & Housing investment elasticity param. & 2.12 (Sommer \& Sullivan 2017)\\
$\lambda$ & Max loan-to-value ratio & 0.8\\
$n$ & Number of housing sizes & 7\\
\bottomrule
\end{tabular}
\end{table}
\end{frame}

\begin{frame}{Internal calibration --- Preferences and housing demand}
\small
\begin{table}
\centering
\begin{tabular}{lll}
\toprule
Parameter & Target moment & Value\\
\midrule
$\beta$ & Net wealth/GDP = 3.25 & 0.944\\
$\omega_{\ell}$ & Avg. time working = 40\% & 8.5\\
$\mathcal{N}$ & Avg. VSL = \$7.5M & 340\\
$\theta_1$ & Avg. rent/income = 0.33 & 0.278\\
$h_1$ & Renters w/ rent/income $\ge 0.5$ = 16.2\% & $0.156\times\bar h_o$\\
$h_6$ & Homeownership rate = 66.1\% & $5.17\times \frac{p^h}{\bar y_{\ell,o}}$\\
$\underline{h}$ & Homelessness rate = 0.30\% & $0.21\times h_1$\\
\bottomrule
\end{tabular}
\end{table}
\end{frame}

\begin{frame}{Internal calibration --- Rental supply, drugs, and policy take-up}
\small
\begin{table}
\centering
\begin{tabular}{lll}
\toprule
Parameter & Target moment & Value\\
\midrule
$\gamma_2$ & Rental supply elasticity = 1.4 & 2.35\\
$\gamma_1$ & Rent/price ratio = 8.3\% & Implied\\
$J_A$ & Avg. age of addicts = 40 & 55\\
$\theta_{0,2}$ & Drugs spending/income (non-addicts) = 0.11 & 0.13\\
$\theta_{1,2}$ & Drugs spending/income (addicts) = 0.45 & 0.45\\
$b$ & Drug users/population = 4\% & $3.7\times 10^{-8}$\\
$\pi_{d,low}$ & Experimenters/addicts = 1 & 0.037\\
$\bar v$ & Light experimenters/experimenters = 0.868 & $5.41\times b$\\
$\kappa_{d,\phi}$ & Participation elasticity = $-0.63$ & 6.9\\
$\pi^{base}_{eu}$ & Unemployment rate = 4.9\% & 0.016\\
$\kappa_{d,ue}$ & Addicts unemployment rate = 26.9\% & 0.37\\
$\kappa_{H,d}$ & Share of homeless who are addicted = 37\% & $\tfrac{1}{6}$\\
$\pi^{V}$ & Voucher recipients/eligible = 26.6\% & 0.06\\
$\kappa_f$ & Food stamps/GDP = 0.42\% & $0.126\times y$\\
$\tau_{\ell}$ & Avg. labor income tax rate = 22.4\% & 0.75\\
\bottomrule
\end{tabular}
\end{table}
\end{frame}

% ---------------- Model performance ----------------
\section{Model performance}

\begin{frame}{Model performance --- Homelessness and addiction}
\small
\begin{table}
\centering
\begin{tabular}{lc}
\toprule
Statistic & Model\\
\midrule
Homeless unemployed (\%) & 100\\
Addicts unemployed (\%) & 22.8\\
Homeless avg. income & $0.12\times$ avg. income\\
Homeless avg. assets & $0.02\times$ avg. assets\\
\bottomrule
\end{tabular}
\end{table}
\end{frame}

\begin{frame}{Model performance --- Drug use}
\small
\begin{table}
\centering
\begin{tabular}{lc}
\toprule
Statistic & Model\\
\midrule
Experimenters' avg. income & $0.2\times$ avg. income\\
Experimenters' avg. assets & $0.15\times$ avg. assets\\
Experimenters while unemployed (\%) & 64.4\\
Experimenters while homeless (\%) & 8.9\\
Experimenters while retired (\%) & 0\\
Voucher recipients who use drugs (\%) & 14.5\\
Voucher recipients who experiment (\%) & 7.9\\
\bottomrule
\end{tabular}
\end{table}
\end{frame}

\begin{frame}{Model performance --- Shares of income and wealth}
\small
\begin{table}
\centering
\begin{tabular}{lcc}
\toprule
Statistic (\%) & Model & Data (SCF)\\
\midrule
\multicolumn{3}{l}{\textbf{Share of total income}}\\
Bottom 1\% & 0 & 0\\
Bottom 1--5\% & 1.1 & 0.3\\
Bottom 5--10\% & 2.6 & 0.6\\
\addlinespace
\multicolumn{3}{l}{\textbf{Share of total net wealth}}\\
Bottom 1\% & 0 & -0.2\\
Bottom 1--5\% & 0 & -0.1\\
Bottom 5--10\% & 0 & 0\\
\bottomrule
\end{tabular}
\end{table}
\end{frame}

\begin{frame}{Model performance --- Drug use by age group}
\centering
\includegraphics[width=0.92\textwidth]{slide_imgs/page_41.png}
\end{frame}

\begin{frame}{Model performance --- Homelessness by age group}
\centering
\includegraphics[width=0.92\textwidth]{slide_imgs/page_42.png}
\end{frame}

% ---------------- Experiments ----------------
\section{Experiments}

\begin{frame}{First experiment's description}
\begin{itemize}
\item Compare \emph{Treatment First} (TF) vs \emph{Housing First} (HF).
\item Vouchers are not given / are taken away if addicted: $\eta=0$ if $d=1$.
\item Four cases:
\begin{itemize}
\item Partial equilibrium (PE): 1 year after implementation.
\item PE: 5 years after implementation.
\item PE: long run.
\item General equilibrium (GE): prices, bequests, and vouchers adjust.
\end{itemize}
\end{itemize}
\end{frame}

\begin{frame}{Results of ``Treatment First'' experiments}
\small
\begin{table}
\centering
\begin{tabular}{lcccc}
\toprule
Statistic (\%\,$\Delta$) & 1 period, PE & 5 periods, PE & Long run, PE & Long run, GE\\
\midrule
Drug users & -1.59 & -1.77 & -2.39 & -2.36\\
Experimenters & -3.01 & -3.15 & -3.21 & -2.98\\
Addicts & -0.16 & -0.40 & -1.57 & -1.73\\
Homelessness & 15.45 & 17.22 & 17.61 & 17.63\\
\# homeless and addicted & 42.91 & 50.74 & 52.21 & 52.27\\
\# homeless 2 years+ & 2.27 & 21.48 & 22.40 & 22.35\\
Life expectancy (years) & 0.00 & 0.00 & 0.01 & 0.01\\
Welfare newborns & Down & Down & Down & Down\\
Welfare everyone & Down & Down & Down & Down\\
\bottomrule
\end{tabular}
\end{table}
\end{frame}

\begin{frame}{Experiments with high moral hazard}
\begin{itemize}
\item Moral hazard is low in benchmark because voucher probability if homeless is only 5.5\%.
\item Two additional long-run GE experiments (full government budget clears):
\begin{itemize}
\item Experiment 2: Vouchers for everyone eligible (VFE) under HF.
\item Experiment 3: Same as Experiment 2 but under TF (sobriety requirements).
\end{itemize}
\end{itemize}
\end{frame}

\begin{frame}{Results of Experiment 2: VFE with HF (relative to benchmark)}
\small
\begin{table}
\centering
\begin{tabular}{lc}
\toprule
Statistic (\%\,$\Delta$) & Long run, GE\\
\midrule
Rent & 2.45\\
House price & -0.87\\
Average income tax & 0.53\\
Drug users & 36.48\\
Experimenters & 36.18\\
Addicts & 36.79\\
Homelessness & -95.25\\
\# homeless and addicted & -97.43\\
\# homeless 2 years+ & -100.00\\
Life expectancy (years) & -0.21\\
Welfare newborns & Up\\
Welfare everyone & Down\\
\bottomrule
\end{tabular}
\end{table}
\end{frame}

\begin{frame}{Results of Experiment 3: VFE with TF (relative to VFE with HF)}
\small
\begin{table}
\centering
\begin{tabular}{lc}
\toprule
Statistic (\%\,$\Delta$) & Long run, GE\\
\midrule
Drug users & -32.91\\
Experimenters & -33.98\\
Addicts & -31.85\\
Homelessness & 980.27\\
\# homeless and addicted & 4{,}553.93\\
\# homeless 2 years+ & 71.85 (p.p.)\\
Life expectancy (years) & 0.26\\
Welfare newborns & Down\\
Welfare everyone & Up\\
\bottomrule
\end{tabular}
\end{table}
\end{frame}

% ---------------- Conclusion ----------------
\section{Conclusion}

\begin{frame}{Conclusion (part 1)}
\begin{itemize}
\item Homelessness and addiction are pressing U.S. issues affecting millions.
\item Many homeless also suffer from addiction.
\item Interconnectedness makes policy design hard:
\begin{itemize}
\item Prioritizing immediate housing can help short term,
\item but may create moral hazard if individuals do not internalize the costs of addiction.
\end{itemize}
\item HF camp: immediate help is more effective despite moral hazard.
\item TF camp: homelessness rises because of moral hazard.
\end{itemize}
\end{frame}

\begin{frame}{Conclusion (part 2)}
\begin{itemize}
\item Build a GE model with rational drug use and endogenous homelessness.
\item Calibrate to key U.S. data on homelessness and addiction.
\item Evaluate short- and long-run implications of TF vs HF.
\item When moral hazard is low (benchmark): TF leads to short-run pain for little long-run gain.
\item When moral hazard is high (if vouchers are universal): TF reduces addiction substantially,
but homelessness still rises significantly.
\end{itemize}
\end{frame}

\end{document}
